% New deduction in run JAC1964-20261121, independently source-checked.
% /root/programme_audit; GPT-6 Astra Pro; effort setting not exposed.

\subsection*{JAC-P09. A geometric generic axis has no variable algebraic classes}

\begin{proposition}[JAC-P09]
\label{jac:P09}
Let $X\subset\mathbf P^N_k$ be an abelian variety of dimension
$2m\ge4$ over an algebraically closed field $k$ of arbitrary
characteristic. Fix $d\ge3$. Let $(H_1,H_2)$ be a geometric generic
ordered pair of degree-$d$ hypersurfaces, and put
\[
 Y_1=X\cap H_1,\qquad B=X\cap H_1\cap H_2,
 \qquad \kappa:B\hookrightarrow X.
\]
All varieties hereafter are taken over the algebraic closure $k'$ of
the function field of this parameter space. Then for every invertible
prime $\ell$,
\[
 \ker\!\left(\kappa_*:\operatorname{Alg}^{m-1}(B)
             \longrightarrow H^{2m+2}(X,\mathbf Q_\ell(m+1))\right)=0.
                                                        \tag{JAC.P09.1}
\]
Equivalently, the variable middle cohomology of this geometric generic
axis contains no nonzero rational algebraic cycle class.

The pencil generated by $(H_1,H_2)$ is Lefschetz. Accepting 032,
Corollary~8.4, its signed geometric height on the entire group
$\operatorname{CH}^{m}(Y/K)_{\mathrm{hom},\mathbf Q}$, where
$K=k'(\mathbf P^1)$ and $Y/K$ is its generic fibre, is positive
semidefinite. Its radical is exactly
\[
 \operatorname{im}\!\left(
  \operatorname{CH}^{m}(\operatorname{Bl}_B X)_{
       \mathrm{hom},\mathbf Q}
     \longrightarrow\operatorname{CH}^{m}(Y)_{\mathbf Q}\right),
                                                        \tag{JAC.P09.2}
\]
and the resulting quotient, with its height form, is canonically
isometric to
\[
 \left(\ker\!\left(L_h:\operatorname{Alg}^{m}(X)
                        \longrightarrow H^{2m+2}(X)(m+1)\right),
       \ (z,z')\longmapsto(-1)^m\deg_X(zz')\right),
                                                        \tag{JAC.P09.3}
\]
where $h$ is the degree-$d$ section class.
\end{proposition}

\begin{proof}
We first identify the parameter field precisely. Choose generic
coefficients for $H_1$, let $k_1$ be an algebraic closure of their
function field inside $k'$, and regard $Y_1$ as a smooth projective
variety over $k_1$, with the projective embedding induced from
$X\subset\mathbf P^N$. The independent coefficients of $H_2$ are
still generic over $k_1$. Thus $B/k'$ is the geometric generic member
of the universal family of smooth degree-$d$ hypersurface sections
of this fixed embedded $Y_1/k_1$. This is a statement about its
geometric generic fibre, not about every closed parameter point in
a Zariski open subset.

The dimension of $Y_1$ is $2m-1\ge3$, which is odd.
Katz--Pandharipande, \emph{Inequalities related to Lefschetz pencils
and integrals of Chern classes}, Corollaries~2 and~3, author PDF
p.~3, show that its variable middle local system has rank $N_d>8$
and geometric monodromy group the full orthogonal group $O(N_d)$.
The only exceptional case in those corollaries is a linearly
embedded $\mathbf P^3$ with $d=3$. It cannot occur here:
weak Lefschetz gives
$H^1(Y_1,\mathbf Q_\ell)\simeq H^1(X,\mathbf Q_\ell)\ne0$,
whereas $H^1(\mathbf P^3,\mathbf Q_\ell)=0$.
The hypotheses in the paper's introduction are only that $k_1$ is
algebraically closed and $\ell$ is invertible. Corollary~3 imposes
no restriction on the characteristic. In particular this use of
orthogonal monodromy is valid in characteristic two.

Write $i:B\hookrightarrow Y_1$ and $j_1:Y_1\hookrightarrow X$.
Hard Lefschetz on $Y_1$ makes the ambient subspace
$i^*H^{2m-2}(Y_1)$ nondegenerate for the middle pairing on $B$.
Consequently its orthogonal complement
\[
 E_B=\ker\!\left(i_*:H^{2m-2}(B)
                         \longrightarrow H^{2m}(Y_1)(1)\right)
\]
identifies equivariantly with the variable quotient local system
used by Katz--Pandharipande. The projection onto $E_B$ commutes with
geometric monodromy; it need not be asserted algebraic for this step.

Every algebraic cycle on the geometric generic fibre is defined
over a finite algebraic extension of the generic parameter field.
Its cohomology class therefore has a finite orbit under the
geometric fundamental group. Purely inseparable field extensions
do not change this assertion: their absolute Galois groups and
prime-to-characteristic \^{e}tale cohomology are unchanged, and the
separable part of the field of definition has finite degree.
The variable projection of the class also has a finite orbit.
Its stabilizer has finite index, so the Zariski closure of that
stabilizer contains the identity component $SO(N_d)$ of $O(N_d)$.
The standard representation of $SO(N_d)$, with $N_d>8$, has no
nonzero fixed vector. The variable projection is therefore zero.

Finally,
\[
 (j_1)_*:H^{2m}(Y_1)\longrightarrow H^{2m+2}(X)(1)
\]
is an isomorphism, by Poincar\'e duality and weak Lefschetz in degree
$2m-2$. Since $\kappa_*=(j_1)_*i_*$, its kernel on
$H^{2m-2}(B)$ is precisely $E_B$. The preceding paragraph proves
(JAC.P09.1), including for rational algebraic combinations of cycles.

A geometric generic ordered independent pair determines a geometric
generic line in the projective parameter space of degree-$d$
sections. SGA~7~II, XVII, Theorem~2.5.1, printed p.~217, ensures
that this is a Lefschetz pencil for $d\ge2$. Its axis is smooth.
JAC-P07 supplies condition~A, so H9-T7 applies to all homologically
trivial cycles over its original pencil field $K$. In the
orthogonal decomposition of H9-T7, formula (JAC.P09.1) says that
the axis summand $A_B$ is zero. Only the ambient primitive summand
remains. Its form is positive definite by 032, Corollary~8.4;
H9-T7 then gives (JAC.P09.2)--(JAC.P09.3) and the asserted
semipositivity.
\end{proof}

\paragraph{Scope of the genericity and of the height assertion.}
The assertion about $B$ is geometric generic. It does not assert
vanishing of $A_B$ for every sufficiently general closed point over
the original field, nor for a fixed special axis. One may equivalently
describe the pencil as the geometric generic pencil in the
Grassmannian, after adjoining a generic ordered frame and taking an
algebraic closure. The result concerns cycles defined over the
original pencil function field $k'(\mathbf P^1)$; it makes no claim
after an arbitrary finite covering of the pencil base, whose total
space need not be a blowup of $X$.

\paragraph{Primary source for the new monodromy input.}
The author-hosted PDF is
\url{https://web.math.princeton.edu/~nmk/Inequalities15.pdf}.
Its introduction, p.~1, defines the variable quotient and geometric
monodromy over an algebraically closed field with one invertible
prime. Theorem~0, p.~2, cites Katz's theorem that $N_d>8$ and
$d\ge3$ force full orthogonal monodromy in even fibre dimension.
Corollaries~2--3, p.~3, remove the rank hypothesis with the unique
linear cubic-surface exception. Those pages were inspected in the
original author PDF in this run. The characteristic restriction
appearing later in the separate Theorem~4 is not a hypothesis of
Corollary~3.
